\section{Measure the scan and loop separately}\label{packet-measurement}
The two verified implementations make a useful comparison. The moving-pointer
version shown in the increments scans the eight header bytes again to compute
their checksum. The indexed version accumulates those bytes while decoding them.
Its body loop has a top test, indexed address calculation, and unconditional
backward branch. Eliminating eight repeated reads therefore changes both fixed
setup work and the per-payload-byte instruction cost.

\Needspace{10\baselineskip}
\begin{longtable}{rrrrr}
\toprule
Payload bytes&Pointer instr.&Indexed instr.&Pointer reads&Indexed reads\\\midrule
0&111&70&17&9\\
3&132&94&20&12\\
64&559&582&81&73\\\bottomrule
\end{longtable}
These are native forward-execution counts for one successful independently called
function in the default memory layout. Counting begins immediately before the
parser executes and stops before its supplied return continuation. Argument setup,
assembly, session construction, main, and the exit syscall are excluded. Input
reads count simulator data-access notices overlapping the arena; each read in these
two implementations is a byte load. The generated valid packets use sequence 4660, flags one, and payload
byte \(i\) equal to \((129i+127)\bmod256\). Their three-byte body is
\texttt{7f 00 81}; it differs from the worked packet while following the same
successful control path for that length.

The indexed version executes fewer instructions for the empty and short examples,
but more for the 64-byte payload. For these implementations the measured successful
paths follow \(111+7L\) and \(70+8L\) instructions. The fixed saving is exhausted
at 41 payload bytes. Both loops are correct; the lower load count alone does not
predict the lower instruction count for every valid input.

Instruction counts are not host elapsed time or a target processor cycle model.
A different loop, register allocation, calling convention, or memory configuration
can change them. Use these measurements to locate repeated work and formulate an
optimization, then rerun the behavioral and read-boundary checks. The counts are
an illustration, never an exact-count requirement imposed on student submissions.
